\textbf{1. Kiến thức cần nhớ}
\begin{itemize}
\item \textbf{Phương trình bậc hai một ẩn} là phương trình có dạng:
$$\boxed{ax^2 + bx + c = 0\quad (a \ne 0)}$$
\item \textbf{Công thức nghiệm}: Biệt thức $\Delta = b^2 - 4ac$:
\begin{itemize}
\item Nếu $\Delta > 0$: Phương trình có \textbf{hai nghiệm phân biệt}:
$$\boxed{x_1 = \dfrac{-b + \sqrt{\Delta}}{2a};\quad x_2 = \dfrac{-b - \sqrt{\Delta}}{2a}}$$
\item Nếu $\Delta = 0$: Phương trình có \textbf{nghiệm kép}:
$$\boxed{x_1 = x_2 = -\dfrac{b}{2a}}$$
\item Nếu $\Delta < 0$: Phương trình \textbf{vô nghiệm}.
\end{itemize}
\item \textbf{Công thức nghiệm thu gọn} (khi $b = 2b'$, $\Delta' = b'^2 - ac$):
\begin{itemize}
\item $\Delta' > 0 \implies x_{1,2} = \dfrac{-b' \pm \sqrt{\Delta'}}{a}$.
\item $\Delta' = 0 \implies x_1 = x_2 = -\dfrac{b'}{a}$.
\item $\Delta' < 0 \implies$ vô nghiệm.
\end{itemize}
\end{itemize}

\textbf{2. Phương pháp giải các dạng}

%% DANG 9D6?19-1
\textbf{Dạng 1. Giải phương trình bậc hai khuyết}
\begin{enumerate}
\item Khuyết $c$ ($ax^2 + bx = 0$): Đặt nhân tử chung $x(ax + b) = 0 \iff x = 0$ hoặc $x = -\dfrac{b}{a}$.
\item Khuyết $b$ ($ax^2 + c = 0$): $x^2 = -\dfrac{c}{a}$. Có nghiệm $x = \pm\sqrt{-\dfrac{c}{a}}$ khi $ac < 0$; vô nghiệm khi $ac > 0$.
\end{enumerate}

%% DANG 9D6?19-2
\textbf{Dạng 2. Giải bằng công thức nghiệm, công thức nghiệm thu gọn}
\begin{enumerate}
\item Xác định các hệ số $a, b, c$.
\item Tính $\Delta = b^2 - 4ac$ (hoặc $\Delta' = b'^2 - ac$).
\item Xét dấu của $\Delta$ rồi áp dụng công thức nghiệm tương ứng.
\end{enumerate}

%% DANG 9D6?19-3
\textbf{Dạng 3. Số nghiệm của phương trình bậc hai chứa tham số}
\begin{enumerate}
\item Chú ý hệ số $a$: Nếu $a$ chứa tham số, phải xét trường hợp $a = 0$ trước.
\item Khi $a \ne 0$:
\begin{itemize}
\item Có $2$ nghiệm phân biệt $\iff \Delta > 0$.
\item Có nghiệm kép $\iff \Delta = 0$.
\item Vô nghiệm $\iff \Delta < 0$.
\item Có nghiệm $\iff \Delta \ge 0$.
\end{itemize}
\end{enumerate}
